% compile with XeLaTeX
\documentclass[dvipsnames,mathserif]{beamer}
\setbeamertemplate{footline}[frame number]
\setbeamercolor{footline}{fg=black}
\setbeamerfont{footline}{series=\bfseries}
\usepackage{tikz}
\usepackage{xcolor}
\usepackage{booktabs}

\usetheme{Darmstadt}

\makeatletter
\newcommand{\RTListe}{\raggedleft\rightskip\leftm}
\newcommand{\leftm}{\@totalleftmargin}
\makeatother

\setbeamertemplate{frametitle}
{\vspace*{-1mm}
  \nointerlineskip
    \begin{beamercolorbox}[sep=0.3cm,ht=2.2em,wd=\paperwidth]{frametitle}
        \vbox{}\vskip-2ex%
        \strut\hskip1ex\insertframetitle\strut
        \vskip-0.8ex%
    \end{beamercolorbox}
}
\makeatletter
\setbeamertemplate{subsection in toc}
{\leavevmode\rightskip=5ex%
  \llap{\raise0.1ex\beamer@usesphere{subsection number projected}{bigsphere}\kern1ex}%
  \inserttocsubsection\par%
}
\makeatother

\setbeamertemplate{itemize item}{\scriptsize\raise1.25pt\hbox{\donotcoloroutermaths$\blacktriangleleft$}} 

\defbeamertemplate{enumerate item}{square2}
{\LR{
    \hbox{%
    \usebeamerfont*{item projected}%
    \usebeamercolor[bg]{item projected}%
    \vrule width2.25ex height1.85ex depth.4ex%
    \hskip-2.25ex%
    \hbox to2.25ex{%
      \hfil%
      {\color{fg}\insertenumlabel}%
      \hfil}%
  }%
}}

\setbeamertemplate{enumerate item}[square2]
\setbeamertemplate{navigation symbols}{}

% FOSSEE LOGO ON EVERY PAGE
\logo{\includegraphics[height=0.9cm]{fossee_logo.png}}

\setbeamertemplate{caption}[numbered]
\begin{document}

\rightskip\rightmargin
\title{\textbf{Simulation Orchestration \& Async Integration} \\ Technical Progress Report (Phase 3)}
\author{\Large \textbf{Pranjal Kumar Shukla}}
\institute{\large\textbf{FOSSEE Internship -- Progress Report}\\[4pt]
SCAI School\\[2pt]
VIT Bhopal University}
\footnotesize{\date{\today}}


\begin{frame}
\maketitle
\end{frame}

%% ─────────────────────────────────────────
%% 1. ORCHESTRATION & ASYNC
%% ─────────────────────────────────────────
\section{Solver Automation}
\begin{frame}{1. Multi-Solver Integration}
    \begin{itemize}
        \item \textbf{Objective}: Evolve the addon from a pure meshing utility into a complete CFD orchestrator.
        \item \textbf{Implementation}: Integrated GUI controls to configure real-world physical bounds.
        \begin{itemize}
            \item Selection between \texttt{icoFoam} (Transient) and \texttt{simpleFoam} (Steady-State).
            \item Parametric inputs for Kinematic Viscosity ($\nu$) and Inlet Velocity Vectors.
        \end{itemize}
        \item \textbf{Dynamic Dictionary Gen}: The backend Python engine cleanly writes the \texttt{0/U}, \texttt{0/p}, and \texttt{system/controlDict} configurations natively across OS boundaries.
    \end{itemize}
    \begin{center}
        \includegraphics[width=0.7\textwidth]{ui_week3_panels.png}
        \caption*{Revamped Sub-Panel UI for Mesh, Physics, and Post-Processing}
    \end{center}
\end{frame}

\begin{frame}{2. Asynchronous Execution Pipeline}
    \begin{itemize}
        \item \textbf{The Bottleneck}: OpenFOAM solver iterations freeze Blender's main UI thread, forcing users to "blindly wait."
        \item \textbf{The Solution}: 
        \begin{itemize}
            \item Implemented Python \texttt{threading.Thread} for all terminal background processes.
            \item Registered a recurring \texttt{bpy.app.timers} hook to constantly poll Global State variables.
        \end{itemize}
        \item \textbf{Result}: The Blender viewport remains 100\% fluid and functional while \texttt{snappyHexMesh} or \texttt{simpleFoam} chew through CPU cycles.
    \end{itemize}
    \begin{center}
        \includegraphics[width=0.6\textwidth]{async_status.png}
        \caption*{Real-Time Non-Blocking Status Indicator}
    \end{center}
\end{frame}

%% ─────────────────────────────────────────
%% 2. BUG FIXES & OF-11 QUIRKS
%% ─────────────────────────────────────────
\section{Engineering Challenges}
\begin{frame}{3. Resolving Python Architecture Bugs}
    Developing highly compartmentalized Blender extensions led to severe edge cases:
    \begin{itemize}
        \item \textbf{\texttt{bl\_info} Registration}: A typo in the registration header (\texttt{maxbl\_info}) caused the entire Addon to become invisible to Blender's plugin parser.
        \item \textbf{Module NotFound Errors}: Attempting standard Python \texttt{import run\_cfmesh} failed catastrophically inside Blender's complex system path. 
        \item \textbf{The Fix}: Implemented strict **Relative Imports** (\texttt{from . import run\_cfmesh}) forcing Blender to reliably resolve the module within its own localized ZIP container.
    \end{itemize}
\end{frame}

\begin{frame}{4. Navigating OpenFOAM 11's New Paradigm}
    \begin{alertblock}{The "simpleFoam is Superseded" Error}
        OpenFOAM 11 entirely deleted standard standalone solver aliases in favor of a universal module execution layer.
    \end{alertblock}
    \begin{itemize}
        \item \textbf{The Migration}: We transitioned the native Blender call from \texttt{simpleFoam} to \texttt{foamRun -solver incompressibleFluid}.
        \item \textbf{Dictionary Overhauls}: 
        \begin{itemize}
            \item OF-11 dynamically determines if a simulation is Steady vs Transient based on the \texttt{ddtSchemes} variable.
            \item Implemented logic to inject \texttt{steadyState} and \texttt{SIMPLE} equation blocks dynamically depending on the user's dropdown choice in the Blender UI.
            \item Added explicit generation for \texttt{constant/momentumTransport} to satisfy OF-11's modular turbulence checking.
        \end{itemize}
    \end{itemize}
\end{frame}

\begin{frame}{5. Fixing the Hollow Mesh Failure}
    \begin{itemize}
        \item \textbf{Issue}: \texttt{snappyHexMesh} frequently returned an un-meshed solid cube instead of the desired object.
        \item \textbf{Cause \#1}: The \texttt{locationInMesh} point was hardcoded to \texttt{(0,0,0)}. If an object sits at the exact origin, OpenFOAM attempts to mesh the \textit{inside} of it instead of the surrounding wind tunnel!
        \item \textbf{Cause \#2}: User-favorites like the "Suzanne" monkey are not watertight surfaces.
        \item \textbf{Resolution}: Shifted the internal meshing coordinate to \texttt{(1.9, 1.9, 1.9)} to guarantee location in the external fluid domain.
    \end{itemize}
\end{frame}

%% ─────────────────────────────────────────
%% 3. VISUALIZATION & WRAP UP
%% ─────────────────────────────────────────
\section{Visualization}
\begin{frame}{6. ParaView Integration}
    \begin{columns}
    \column{0.5\textwidth}
    The ultimate goal of Phase 3 was seamless End-to-End orchestration.
    \begin{itemize}
        \item A new native UI button generates the `.foam` environment trigger file and launches ParaView automatically.
        \item Fully bridges the gap between Blender's 3D workspace, OpenFOAM's simulation engine, and ParaView's post-processing rendering array.
    \end{itemize}
    \column{0.5\textwidth}
    \begin{center}
        \includegraphics[width=0.9\textwidth]{paraview_result.png}
        \caption*{Steady-State Pressure visualized in ParaView}
    \end{center}
    \end{columns}
\end{frame}

\section{Next Phase}
\begin{frame}{7. Toward Week 4: Advanced Scenarios}
    \begin{enumerate}
        \item \textbf{Dynamic Meshing}: Implementing parameter bounds for moving objects.
        \item \textbf{In-Blender Post Processing}: Basic pseudo-color visualization onto Blender surfaces without needing ParaView.
        \item \textbf{Turbulence Modeling}: Expanding the solver boundaries to support sophisticated \texttt{k-epsilon} and \texttt{k-omega} RANS models.
    \end{enumerate}
\end{frame}

\begin{frame}
    \begin{center}
        \Huge \textbf{Thank You!} \\
        \vspace{0.5cm}
        \large Questions?
    \end{center}
\end{frame}

\end{document}
